\IfFileExists{stacks-project.cls}{%
\documentclass{stacks-project}
}{%
\documentclass{amsart}
}

% For dealing with references we use the comment environment
\usepackage{verbatim}
\newenvironment{reference}{\comment}{\endcomment}
%\newenvironment{reference}{}{}
\newenvironment{slogan}{\comment}{\endcomment}
\newenvironment{history}{\comment}{\endcomment}

% For commutative diagrams we use Xy-pic
\usepackage[all]{xy}

% We use 2cell for 2-commutative diagrams.
\xyoption{2cell}
\UseAllTwocells

% We use multicol for the list of chapters between chapters
\usepackage{multicol}

% This is generally recommended for better output
\usepackage{lmodern}
\usepackage[T1]{fontenc}

% For cross-file-references
\usepackage{xr-hyper}

% Package for hypertext links:
\usepackage{hyperref}

% For any local file, say "hello.tex" you want to link to please
% use \externaldocument[hello-]{hello}
\externaldocument[introduction-]{introduction}
\externaldocument[conventions-]{conventions}
\externaldocument[sets-]{sets}
\externaldocument[categories-]{categories}
\externaldocument[topology-]{topology}
\externaldocument[sheaves-]{sheaves}
\externaldocument[sites-]{sites}
\externaldocument[stacks-]{stacks}
\externaldocument[fields-]{fields}
\externaldocument[algebra-]{algebra}
\externaldocument[brauer-]{brauer}
\externaldocument[homology-]{homology}
\externaldocument[derived-]{derived}
\externaldocument[simplicial-]{simplicial}
\externaldocument[more-algebra-]{more-algebra}
\externaldocument[smoothing-]{smoothing}
\externaldocument[modules-]{modules}
\externaldocument[sites-modules-]{sites-modules}
\externaldocument[injectives-]{injectives}
\externaldocument[cohomology-]{cohomology}
\externaldocument[sites-cohomology-]{sites-cohomology}
\externaldocument[dga-]{dga}
\externaldocument[dpa-]{dpa}
\externaldocument[sdga-]{sdga}
\externaldocument[hypercovering-]{hypercovering}
\externaldocument[schemes-]{schemes}
\externaldocument[constructions-]{constructions}
\externaldocument[properties-]{properties}
\externaldocument[morphisms-]{morphisms}
\externaldocument[coherent-]{coherent}
\externaldocument[divisors-]{divisors}
\externaldocument[limits-]{limits}
\externaldocument[varieties-]{varieties}
\externaldocument[topologies-]{topologies}
\externaldocument[descent-]{descent}
\externaldocument[perfect-]{perfect}
\externaldocument[more-morphisms-]{more-morphisms}
\externaldocument[flat-]{flat}
\externaldocument[groupoids-]{groupoids}
\externaldocument[more-groupoids-]{more-groupoids}
\externaldocument[etale-]{etale}
\externaldocument[chow-]{chow}
\externaldocument[intersection-]{intersection}
\externaldocument[pic-]{pic}
\externaldocument[weil-]{weil}
\externaldocument[adequate-]{adequate}
\externaldocument[dualizing-]{dualizing}
\externaldocument[duality-]{duality}
\externaldocument[discriminant-]{discriminant}
\externaldocument[derham-]{derham}
\externaldocument[local-cohomology-]{local-cohomology}
\externaldocument[algebraization-]{algebraization}
\externaldocument[curves-]{curves}
\externaldocument[resolve-]{resolve}
\externaldocument[models-]{models}
\externaldocument[functors-]{functors}
\externaldocument[equiv-]{equiv}
\externaldocument[pione-]{pione}
\externaldocument[etale-cohomology-]{etale-cohomology}
\externaldocument[proetale-]{proetale}
\externaldocument[relative-cycles-]{relative-cycles}
\externaldocument[more-etale-]{more-etale}
\externaldocument[trace-]{trace}
\externaldocument[crystalline-]{crystalline}
\externaldocument[spaces-]{spaces}
\externaldocument[spaces-properties-]{spaces-properties}
\externaldocument[spaces-morphisms-]{spaces-morphisms}
\externaldocument[decent-spaces-]{decent-spaces}
\externaldocument[spaces-cohomology-]{spaces-cohomology}
\externaldocument[spaces-limits-]{spaces-limits}
\externaldocument[spaces-divisors-]{spaces-divisors}
\externaldocument[spaces-over-fields-]{spaces-over-fields}
\externaldocument[spaces-topologies-]{spaces-topologies}
\externaldocument[spaces-descent-]{spaces-descent}
\externaldocument[spaces-perfect-]{spaces-perfect}
\externaldocument[spaces-more-morphisms-]{spaces-more-morphisms}
\externaldocument[spaces-flat-]{spaces-flat}
\externaldocument[spaces-groupoids-]{spaces-groupoids}
\externaldocument[spaces-more-groupoids-]{spaces-more-groupoids}
\externaldocument[bootstrap-]{bootstrap}
\externaldocument[spaces-pushouts-]{spaces-pushouts}
\externaldocument[spaces-chow-]{spaces-chow}
\externaldocument[groupoids-quotients-]{groupoids-quotients}
\externaldocument[spaces-more-cohomology-]{spaces-more-cohomology}
\externaldocument[spaces-simplicial-]{spaces-simplicial}
\externaldocument[spaces-duality-]{spaces-duality}
\externaldocument[formal-spaces-]{formal-spaces}
\externaldocument[restricted-]{restricted}
\externaldocument[spaces-resolve-]{spaces-resolve}
\externaldocument[formal-defos-]{formal-defos}
\externaldocument[defos-]{defos}
\externaldocument[cotangent-]{cotangent}
\externaldocument[examples-defos-]{examples-defos}
\externaldocument[algebraic-]{algebraic}
\externaldocument[examples-stacks-]{examples-stacks}
\externaldocument[stacks-sheaves-]{stacks-sheaves}
\externaldocument[criteria-]{criteria}
\externaldocument[artin-]{artin}
\externaldocument[quot-]{quot}
\externaldocument[stacks-properties-]{stacks-properties}
\externaldocument[stacks-morphisms-]{stacks-morphisms}
\externaldocument[stacks-limits-]{stacks-limits}
\externaldocument[stacks-cohomology-]{stacks-cohomology}
\externaldocument[stacks-perfect-]{stacks-perfect}
\externaldocument[stacks-introduction-]{stacks-introduction}
\externaldocument[stacks-more-morphisms-]{stacks-more-morphisms}
\externaldocument[stacks-geometry-]{stacks-geometry}
\externaldocument[moduli-]{moduli}
\externaldocument[moduli-curves-]{moduli-curves}
\externaldocument[examples-]{examples}
\externaldocument[exercises-]{exercises}
\externaldocument[guide-]{guide}
\externaldocument[desirables-]{desirables}
\externaldocument[coding-]{coding}
\externaldocument[obsolete-]{obsolete}
\externaldocument[fdl-]{fdl}
\externaldocument[index-]{index}

% Theorem environments.
%
\theoremstyle{plain}
\newtheorem{theorem}[subsection]{Theorem}
\newtheorem{proposition}[subsection]{Proposition}
\newtheorem{lemma}[subsection]{Lemma}

\theoremstyle{definition}
\newtheorem{definition}[subsection]{Definition}
\newtheorem{example}[subsection]{Example}
\newtheorem{exercise}[subsection]{Exercise}
\newtheorem{situation}[subsection]{Situation}

\theoremstyle{remark}
\newtheorem{remark}[subsection]{Remark}
\newtheorem{remarks}[subsection]{Remarks}

\numberwithin{equation}{subsection}

% Macros
%
\def\lim{\mathop{\mathrm{lim}}\nolimits}
\def\colim{\mathop{\mathrm{colim}}\nolimits}
\def\Spec{\mathop{\mathrm{Spec}}}
\def\Hom{\mathop{\mathrm{Hom}}\nolimits}
\def\Ext{\mathop{\mathrm{Ext}}\nolimits}
\def\SheafHom{\mathop{\mathcal{H}\!\mathit{om}}\nolimits}
\def\SheafExt{\mathop{\mathcal{E}\!\mathit{xt}}\nolimits}
\def\Sch{\mathit{Sch}}
\def\Mor{\mathop{\mathrm{Mor}}\nolimits}
\def\Ob{\mathop{\mathrm{Ob}}\nolimits}
\def\Sh{\mathop{\mathit{Sh}}\nolimits}
\def\NL{\mathop{N\!L}\nolimits}
\def\CH{\mathop{\mathrm{CH}}\nolimits}
\def\proetale{{pro\text{-}\acute{e}tale}}
\def\etale{{\acute{e}tale}}
\def\QCoh{\mathit{QCoh}}
\def\Ker{\mathop{\mathrm{Ker}}}
\def\Im{\mathop{\mathrm{Im}}}
\def\Coker{\mathop{\mathrm{Coker}}}
\def\Coim{\mathop{\mathrm{Coim}}}

% Boxtimes
%
\DeclareMathSymbol{\boxtimes}{\mathbin}{AMSa}{"02}

%
% Macros for moduli stacks/spaces
%
\def\QCohstack{\mathcal{QC}\!\mathit{oh}}
\def\Cohstack{\mathcal{C}\!\mathit{oh}}
\def\Spacesstack{\mathcal{S}\!\mathit{paces}}
\def\Quotfunctor{\mathrm{Quot}}
\def\Hilbfunctor{\mathrm{Hilb}}
\def\Curvesstack{\mathcal{C}\!\mathit{urves}}
\def\Polarizedstack{\mathcal{P}\!\mathit{olarized}}
\def\Complexesstack{\mathcal{C}\!\mathit{omplexes}}
% \Pic is the operator that assigns to X its picard group, usage \Pic(X)
% \Picardstack_{X/B} denotes the Picard stack of X over B
% \Picardfunctor_{X/B} denotes the Picard functor of X over B
\def\Pic{\mathop{\mathrm{Pic}}\nolimits}
\def\Picardstack{\mathcal{P}\!\mathit{ic}}
\def\Picardfunctor{\mathrm{Pic}}
\def\Deformationcategory{\mathcal{D}\!\mathit{ef}}
\usepackage{cleveref}
\crefname{originalreference}{original reference}{original references}
\Crefname{originalreference}{Original reference}{Original references}
\makeatletter
\newlabel{definition-homotopy}{{019K}{source}{}{sourcepointer-1}{}}
\newlabel{definition-homotopy@cref}{{[originalreference][0][]019K}{[1][0][]source}{}{}{}}
\makeatother
\crefname{originalequation}{original equation}{original equations}
\Crefname{originalequation}{Original equation}{Original equations}
\makeatletter
\newlabel{equation-homotopy-to-homotopy}{{019R}{source}{}{sourcepointer-2}{}}
\newlabel{equation-homotopy-to-homotopy@cref}{{[originalequation][0][]019R}{[1][0][]source}{}{}{}}
\makeatother
\crefname{originallemma}{original lemma}{original lemmas}
\Crefname{originallemma}{Original lemma}{Original lemmas}
\makeatletter
\newlabel{homology-lemma-ses-termwise-split}{{011D}{source}{}{sourcepointer-3}{}}
\newlabel{homology-lemma-ses-termwise-split@cref}{{[originallemma][0][]011D}{[1][0][]source}{}{}{}}
\makeatother
\crefname{originalreference}{original reference}{original references}
\Crefname{originalreference}{Original reference}{Original references}
\makeatletter
\newlabel{lemma-homotopy-s-N}{{019S}{source}{}{sourcepointer-4}{}}
\newlabel{lemma-homotopy-s-N@cref}{{[originalreference][0][]019S}{[1][0][]source}{}{}{}}
\makeatother
\crefname{originalreference}{original reference}{original references}
\Crefname{originalreference}{Original reference}{Original references}
\makeatletter
\newlabel{lemma-map-into-diamond}{{01A3}{source}{}{sourcepointer-5}{}}
\newlabel{lemma-map-into-diamond@cref}{{[originalreference][0][]01A3}{[1][0][]source}{}{}{}}
\makeatother
\crefname{originalreference}{original reference}{original references}
\Crefname{originalreference}{Original reference}{Original references}
\makeatletter
\newlabel{lemma-represent-homotopy}{{01A2}{source}{}{sourcepointer-6}{}}
\newlabel{lemma-represent-homotopy@cref}{{[originalreference][0][]01A2}{[1][0][]source}{}{}{}}
\makeatother
\crefname{originalsection}{original section}{original sections}
\Crefname{originalsection}{Original section}{Original sections}
\makeatletter
\newlabel{section-homotopy}{{019J}{source}{}{sourcepointer-7}{}}
\newlabel{section-homotopy@cref}{{[originalsection][0][]019J}{[1][0][]source}{}{}{}}
\makeatother
\crefname{originalsection}{original section}{original sections}
\Crefname{originalsection}{Original section}{Original sections}
\makeatletter
\newlabel{section-homotopy-abelian}{{019Q}{source}{}{sourcepointer-8}{}}
\newlabel{section-homotopy-abelian@cref}{{[originalsection][0][]019Q}{[1][0][]source}{}{}{}}
\makeatother
\title{Chain homotopies between N(a) and N(b) lift to simplicial homotopies U x Delta[1] -> V}
\author{The Stacks Project Authors}
\date{stacks-project:simplicial.tex:local-snapshot-20261012; GNU FDL 1.2 or later}
\begin{document}
\maketitle
\noindent\textit{Editorial note.} All delivered mathematical text is the author's own, copied from the pinned Stacks Project chapter file sources/stacks-simplicial/pinned-original-source.tex; no statement or proof step was rewritten, reordered or removed.\par\medskip
\noindent\textit{Editorial note.} The counted proof is the author's complete argument: the element-wise computation of delta(sigma) is announced as done 'in case A is the category of abelian groups' (source line 5586) - the standard reduction of morphisms between direct sums to their component maps, which is legitimate because both sides are additive functors in U and the constructed morphisms are natural; every cancellation in the computation is written out and the surviving term is identified; this lemma is the converse companion of the homotopy correspondence of section-homotopy-abelian.\par\medskip
\noindent\textit{Editorial note.} The label lemma-backwards-homotopy names the counted statement, and the proof environment immediately following it in the source is delivered verbatim as its proof.\par\medskip
\noindent\textit{Editorial note.} The counted statement is carried by the author's own lemma environment in the chapter source.\par\medskip
\noindent\textit{Allowed standard prerequisites (editorial).}\begin{itemize}
\item The diamond functor representing homotopy data. Same chapter, \textbackslash{}label\{lemma-represent-homotopy\}. Lines 5522-5526: the triple (N(a), N(b), \{N\_n\}) is represented by a morphism diamond N(U) -> N(V); reused at line 5573 for Psi.
\item The explicit s-level homotopy between e\_0 and e\_1. Same chapter, \textbackslash{}label\{lemma-homotopy-s-N\}. Lines 5547-5552: the explicit homotopy \{h\_n : s(U)\_n -> s(U x Delta[1])\_\{n+1\}\} between s(e\_0) and s(e\_1).
\item The map-into-diamond factorization criterion. Same chapter, \textbackslash{}label\{lemma-map-into-diamond\}. Lines 5553-5556: produces Phi : diamond s(U) -> s(U x Delta[1]) from the homotopy; lines 5576-5581: produces Psi from the factored delta(sigma).
\item Simplicial homotopies and the cylinder object. Same chapter, \textbackslash{}label\{definition-homotopy\} and \textbackslash{}label\{equation-homotopy-to-homotopy\}. Lines 5531-5535 and 5545-5548: e\_0, e\_1 : U -> U x Delta[1] are homotopic with homotopy id; N(h) conventions from section-homotopy-abelian.
\item Termwise split short exact sequences and the connecting morphism. Homology chapter, \textbackslash{}label\{homology-lemma-ses-termwise-split\}. Lines 5584-5588: delta(sigma) is computed as d sigma - sigma d.
\item N is a direct summand of s. Structural property of the normalized complex functor (Same chapter, Section \textbackslash{}ref\{section-homotopy-abelian\}). Lines 5537-5541: the reduction from N to s uses that N is a direct summand of s and that diamond commutes with direct sums.
\end{itemize}\medskip
\ifdefined\hypertarget\hypertarget{sourcepointer-1}{}\fi
\noindent\textit{Source reference.} Original reference 019K: sources/stacks-simplicial/pinned-original-source.tex, source lines 4640--4640. Prerequisite pointer invoked at proof line(s) 5515; it names a standard fact of the same chapter that the counted argument uses, and every step of the counted proof built on it is retained verbatim in the delivered body.\par\medskip
\ifdefined\hypertarget\hypertarget{sourcepointer-2}{}\fi
\noindent\textit{Source reference.} Original equation 019R: sources/stacks-simplicial/pinned-original-source.tex, source lines 4971--4971. Prerequisite pointer invoked at proof line(s) 5532; it names a standard fact of the same chapter that the counted argument uses, and every step of the counted proof built on it is retained verbatim in the delivered body.\par\medskip
\ifdefined\hypertarget\hypertarget{sourcepointer-3}{}\fi
\noindent\textit{Source reference.} Original lemma 011D: sources/stacks-chow/tags.txt, source lines 1359--1359. Prerequisite pointer invoked at proof line(s) 5607; it names a standard fact of a sibling chapter that the counted argument uses, and every step of the counted proof built on it is retained verbatim in the delivered body.\par\medskip
\ifdefined\hypertarget\hypertarget{sourcepointer-4}{}\fi
\noindent\textit{Source reference.} Original reference 019S: sources/stacks-simplicial/pinned-original-source.tex, source lines 5034--5034. Prerequisite pointer invoked at proof line(s) 5533; it names a standard fact of the same chapter that the counted argument uses, and every step of the counted proof built on it is retained verbatim in the delivered body.\par\medskip
\ifdefined\hypertarget\hypertarget{sourcepointer-5}{}\fi
\noindent\textit{Source reference.} Original reference 01A3: sources/stacks-simplicial/pinned-original-source.tex, source lines 5434--5434. Prerequisite pointer invoked at proof line(s) 5559; it names a standard fact of the same chapter that the counted argument uses, and every step of the counted proof built on it is retained verbatim in the delivered body.\par\medskip
\ifdefined\hypertarget\hypertarget{sourcepointer-6}{}\fi
\noindent\textit{Source reference.} Original reference 01A2: sources/stacks-simplicial/pinned-original-source.tex, source lines 5335--5335. Prerequisite pointer invoked at proof line(s) 5523; it names a standard fact of the same chapter that the counted argument uses, and every step of the counted proof built on it is retained verbatim in the delivered body.\par\medskip
\ifdefined\hypertarget\hypertarget{sourcepointer-7}{}\fi
\noindent\textit{Source reference.} Original section 019J: sources/stacks-simplicial/pinned-original-source.tex, source lines 4614--4614. Prerequisite pointer invoked at proof line(s) 5580; it names a standard fact of the same chapter that the counted argument uses, and every step of the counted proof built on it is retained verbatim in the delivered body.\par\medskip
\ifdefined\hypertarget\hypertarget{sourcepointer-8}{}\fi
\noindent\textit{Source reference.} Original section 019Q: sources/stacks-simplicial/pinned-original-source.tex, source lines 4955--4955. Prerequisite pointer invoked at proof line(s) 5518; it names a standard fact of the same chapter that the counted argument uses, and every step of the counted proof built on it is retained verbatim in the delivered body.\par\medskip
\begin{lemma}
\label{lemma-backwards-homotopy}
Let $\mathcal{A}$ be an abelian category.
Let $U$, $V$ be simplicial objects of $\mathcal{A}$.
Let $a, b : U \to V$ be a pair of morphisms.
Assume the corresponding maps of chain complexes
$N(a), N(b) : N(U) \to N(V)$ are homotopic by
a homotopy $\{N_n : N(U)_n \to N(V)_{n + 1}\}$.
Then there exists a homotopy from $a$ to $b$ as in
Definition \ref{definition-homotopy}. Moreover, one can choose the
homotopy $h : U \times \Delta[1] \to V$ such that
$N_n = N(h)_n$ where $N(h)$ is the homotopy coming
from $h$ as in Section \ref{section-homotopy-abelian}.
\end{lemma}

\begin{proof}
Let $(\diamond N(U), \diamond a, \diamond b, \diamond h)$
be as in Lemma \ref{lemma-represent-homotopy} and its proof. By that lemma
there exists a morphism $\diamond N(U) \to N(V)$ representing
the triple $(N(a), N(b), \{N_n\})$. We will show
there exists a morphism
$\psi : N(U \times \Delta[1]) \to \diamond{N(U)}$
such that $\diamond a = \psi \circ N(e_0)$, and
$\diamond b = \psi \circ N(e_1)$. Moreover, we will
show that the homotopy between $N(e_0), N(e_1) :
N(U) \to N(U \times \Delta[1])$ coming from
(\ref{equation-homotopy-to-homotopy}) and
Lemma \ref{lemma-homotopy-s-N} with
$h = \text{id}_{U \times \Delta[1]}$ is mapped via
$\psi$ to the canonical homotopy $\diamond h$ between the two
maps $\diamond a, \diamond b : N(U) \to \diamond{N(U)}$. Certainly this
will imply the lemma.

\medskip\noindent
Note that $N : \text{Simp}(\mathcal{A}) \to \text{Ch}_{\geq 0}(\mathcal{A})$
as a functor is a direct summand of the functor
$s : \text{Simp}(\mathcal{A}) \to \text{Ch}_{\geq 0}(\mathcal{A})$.
Also, the functor $\diamond$ is compatible with direct sums.
Thus it suffices instead to construct a morphism
$\Psi : s(U \times \Delta[1]) \to \diamond{s(U)}$ with the
corresponding properties. This is what we do below.

\medskip\noindent
By Definition \ref{definition-homotopy}
the morphisms $e_0 : U \to U \times \Delta[1]$
and $e_1 : U \to U \times \Delta[1]$ are homotopic
with homotopy $\text{id}_{U \times \Delta[1]}$.
By Lemma \ref{lemma-homotopy-s-N}
we get an explicit homotopy
$\{h_n : s(U)_n \to s(U \times \Delta[1])_{n + 1}\}$
between the morphisms
of chain complexes $s(e_0) : s(U) \to s(U \times \Delta[1])$
and $s(e_1) : s(U) \to s(U \times \Delta[1])$. By
Lemma \ref{lemma-map-into-diamond} above we get a corresponding morphism
$$
\Phi : \diamond{s(U)} \to s(U \times \Delta[1])
$$
According to the construction, $\Phi_n$ restricted to the summand
$s(U)[-1]_n = s(U)_{n - 1}$ of $\diamond{s(U)}_n$
is equal to $h_{n - 1}$. And
$$
h_{n - 1} = \sum\nolimits_{i = 0}^{n - 1}
(-1)^{i + 1} s^n_i \cdot \alpha^n_{i + 1} :
U_{n - 1} \to \bigoplus\nolimits_j U_n \cdot \alpha^n_j.
$$
with obvious notation.

\medskip\noindent
On the other hand, the morphisms $e_i : U \to U \times \Delta[1]$ induce
a morphism $(e_0, e_1) : U \oplus U \to U \times \Delta[1]$.
Denote $W$ the cokernel. Note that, if we write
$(U \times \Delta[1])_n = \bigoplus_{\alpha : [n] \to [1]} U_n \cdot \alpha$,
then we may identify
$W_n = \bigoplus_{i = 1}^n U_n \cdot \alpha^n_i$ with
$\alpha^n_i$ as in Section \ref{section-homotopy}.
We have a commutative diagram
$$
\xymatrix{
0 \ar[r] &
U \oplus U \ar[rd]_{(1, 1)} \ar[r] &
U \times \Delta[1] \ar[d]^\pi \ar[r] &
W \ar[r] &
0 \\
& & U & &
}
$$
This implies we have a similar commutative diagram after
applying the functor $s$. Next, we choose the splittings
$\sigma_n : s(W)_n \to s(U \times \Delta[1])_n$ by mapping
the summand $U_n \cdot \alpha^n_i \subset W_n$ via $(-1, 1)$
to the summands $U_n \cdot \alpha^n_0 \oplus U_n \cdot \alpha^n_i
\subset (U \times \Delta[1])_n$. Note that $s(\pi)_n \circ \sigma_n = 0$.
It follows that $(1, 1) \circ \delta(\sigma)_n = 0$.
Hence $\delta(\sigma)$ factors as in
Lemma \ref{lemma-map-into-diamond}. By that lemma
we obtain a canonical morphism
$\Psi : s(U \times \Delta[1]) \to \diamond{s(U)}$.

\medskip\noindent
To compute $\Psi$ we first compute the morphism
$\delta(\sigma) : s(W) \to s(U)[-1] \oplus s(U)[-1]$.
According to Homology, Lemma \ref{homology-lemma-ses-termwise-split}
and its proof,
to do this we have compute
$$
d_{s(U \times \delta[1]), n} \circ \sigma_n
-
\sigma_{n - 1} \circ d_{s(W), n}
$$
and write it as a morphism into
$U_{n - 1} \cdot \alpha^{n - 1}_0 \oplus U_{n - 1} \cdot \alpha^{n - 1}_n$.
We only do this in case $\mathcal{A}$ is the category of abelian
groups. We use the short hand notation $x_{\alpha}$ for $x \in U_n$
to denote the element $x$ in the summand $U_n \cdot \alpha$
of $(U \times \Delta[1])_n$. Recall that
$$
d_{s(U \times \delta[1]), n}
=
\sum\nolimits_{i = 0}^n (-1)^i d^n_i
$$
where $d^n_i$ maps the summand $U_n \cdot \alpha$
to the summand $U_{n - 1} \cdot (\alpha \circ \delta^n_i)$
via the morphism $d^n_i$ of the simplicial object $U$.
In terms of the notation above this means
$$
d_{s(U \times \delta[1]), n}(x_\alpha) =
\sum\nolimits_{i = 0}^n (-1)^i (d^n_i(x))_{\alpha \circ \delta^n_i}
$$
Starting with $x_\alpha \in W_n$, in other words $\alpha = \alpha^n_j$
for some $j \in \{1, \ldots, n\}$, we see that
$\sigma_n(x_\alpha) = x_\alpha - x_{\alpha^n_0}$ and
hence
$$
(d_{s(U \times \delta[1]), n} \circ \sigma_n)(x_\alpha)
=
\sum\nolimits_{i = 0}^n (-1)^i (d^n_i(x))_{\alpha \circ \delta^n_i}
-
\sum\nolimits_{i = 0}^n (-1)^i (d^n_i(x))_{\alpha^n_0 \circ \delta^n_i}
$$
To compute $d_{s(W), n}(x_\alpha)$, we have to omit
all terms where $\alpha \circ \delta^n_i = \alpha^{n - 1}_0, \alpha^{n - 1}_n$.
Hence we get
$$
\begin{matrix}
(\sigma_{n - 1} \circ d_{s(W), n})(x_\alpha) = \\
\sum\nolimits_{i = 0, \ldots, n\text{ and }
\alpha \circ \delta^n_i \not = \alpha^{n - 1}_0\text{ or }\alpha^{n - 1}_n}
\Big((-1)^i (d^n_i(x))_{\alpha \circ \delta^n_i}
-
(-1)^i (d^n_i(x))_{\alpha^{n - 1}_0}
\Big)
\end{matrix}
$$
Clearly the difference of the two terms is the sum
$$
\sum\nolimits_{i = 0, \ldots, n\text{ and }
\alpha \circ \delta^n_i = \alpha^{n - 1}_0\text{ or }\alpha^{n - 1}_n}
\Big((-1)^i (d^n_i(x))_{\alpha \circ \delta^n_i}
-
(-1)^i (d^n_i(x))_{\alpha^{n - 1}_0}
\Big)
$$
Of course, if $\alpha \circ \delta^n_i = \alpha^{n - 1}_0$
then the term drops out. Recall that $\alpha = \alpha^n_j$
for some $j \in \{1, \ldots, n\}$. The only way
$\alpha^n_j \circ \delta^n_i = \alpha^{n - 1}_n$
is if $j = n$ and $i = n$. Thus we actually get
$0$ unless $j = n$ and in that case we get
$(-1)^n(d^n_n(x))_{\alpha^{n - 1}_n} - (-1)^n(d^n_n(x))_{\alpha^{n - 1}_0}$.
In other words, we conclude the morphism
$$
\delta(\sigma)_n :
W_n
\to
(s(U)[-1] \oplus s(U)[-1])_n = U_{n - 1} \oplus U_{n - 1}
$$
is zero on all summands except $U_n \cdot \alpha^n_n$
and on that summand it is equal to $((-1)^nd^n_n, -(-1)^nd^n_n)$.
(Namely, the first summand of the two corresponds to the factor
with $\alpha^{n - 1}_n$ because that is the map $[n - 1] \to [1]$
which maps everybody to $0$, and hence corresponds to $e_0$.)

\medskip\noindent
We obtain a canonical diagram
$$
\xymatrix{
0 \ar[r] &
s(U) \oplus s(U) \ar[r] \ar[d] &
\diamond{s(U)} \ar[r] \ar[d]^{\Phi}&
s(U)[-1] \ar[r] \ar[d] &
0 \\
0 \ar[r] &
s(U) \oplus s(U) \ar[r] \ar[d] &
s(U \times \Delta[1]) \ar[r] \ar[d]^\Psi &
s(W) \ar[r] \ar[d] &
0 \\
0 \ar[r] &
s(U) \oplus s(U) \ar[r] &
\diamond{s(U)} \ar[r] &
s(U)[-1] \ar[r] &
0
}
$$
We claim that $\Phi \circ \Psi$ is the identity.
To see this it is enough to prove that the composition
of $\Phi$ and $\delta(\sigma)$ as a map
$s(U)[-1] \to s(W) \to s(U)[-1] \oplus s(U)[-1]$ is the
identity in the first factor and minus identity in the second.
By the computations above it is
$((-1)^nd^n_0, -(-1)^nd^n_0) \circ (-1)^n s^n_n = (1, -1)$
as desired.
\end{proof}

\end{document}
